\section{문제 설정과 진단}\label{sec:problem}

\subsection{대상 시스템과 성능 간극}\label{sec:target}

대상은 선행 연구\cite{cha2026fusion}가 구축한 SNN 추론
파이프라인이다. CIFAR-10\cite{krizhevsky2009cifar} 입력을 rate
coding으로 타임스텝마다 이진 스파이크($0$ 또는 $32{,}767$, Q15의
1.0)로 인코딩하고, [Conv $5{\times}5$ + AvgPool + LIF] 두 단과 FC,
출력 LIF를 $T{=}30$회 반복한다. LIF 뉴런은 Q15 고정소수점
감쇠($\beta{=}29{,}491\approx0.9$), 임계값 $32{,}767$, 감산 리셋으로
동작하며, snnTorch\cite{eshraghian2023training}와
Brevitas\cite{brevitas2023}의 QAT 학습
결과\cite{jacob2018quantization}를 정수 전용으로 변환한 것이다. 배포
시 각 Conv+AvgPool 쌍은 $6{\times}6$ stride-2 커널로 대수 퓨전되고
(가중치는 int8 4개의 합인 int16, $|w|\le508$), TFLM 커스텀 연산자가
CMSIS-NN 지원 함수와 함께 이를 실행한다. 표~\ref{tab:model}가 퓨전
후 형상이다. 두 번째 합성곱(CPL2)이 타임스텝당 5.64M MAC의
83.6\%를 차지한다.

\begin{table}[htbp]
\centering
\caption{퓨전 후 레이어 구성과 타임스텝당 MAC.}
\label{tab:model}
\footnotesize
\setlength{\tabcolsep}{3.5pt}
\begin{tabular}{l l l r r}
\hline
레이어 & 입력 & 출력 & $r$(열 길이) & MAC/스텝 \\
\hline
CPL1 & $32{\times}32{\times}3$ & $16{\times}16{\times}32$ & 108 & 884,736 \\
CPL2 & $16{\times}16{\times}32$ & $8{\times}8{\times}64$ & 1,152 & 4,718,592 \\
FC & 4,096 & 10 & -- & 40,960 \\
\hline
\multicolumn{4}{l}{합계 (이미지당 $\times 30$)} & 5,644,288 \\
\hline
\end{tabular}
\begin{flushleft}
\footnotesize 주: $r = k_h k_w C_i$는 합성곱을 행렬곱으로 볼 때의
내적 길이. CPL $=$ Conv+AvgPool+LIF 퓨전 연산자.
\end{flushleft}
\end{table}

선행 실측은 이미지당 $2{,}407.0$ms였다\cite{cha2026fusion}. 이는
비퓨전 표준 경로의 4.0배지만, 유효 처리량으로는 약 70MMAC/s ---
이 칩(Cortex-M7, SMLAD 사이클당 2 MAC, 200MHz에서
400MMAC/s)의 약 18\%다. 남은 5.7배의 간극이 본
연구의 출발점이며, 간극의 원인 후보는 둘이다: 명령을 공급하지
못해서인가(페치 병목), 데이터를 공급하지 못해서인가(로드 병목).
그리고 후자라면, II장의 관찰이 가리키듯 그 데이터의 대부분은 곱해
봐야 0인 피연산자다.

\subsection{진단 방법론: 성능 주장을 인과 주장으로}\label{sec:method}

``최적화 X가 성능을 개선했다''는 주장은 인과 주장이고, 인과 주장이
성립하려면 세 요소가 필요하다. (i) \textbf{조건의 증명}: 라벨이
아니라 실리콘이 답한 상태로 측정 조건을 기록해야 한다. 우리는 매
측정에 \texttt{FLASH\_ACR}$\cdot$\texttt{SCB\_CCR} 레지스터 값과
ITCM 적재 바이트 수를 동봉해, 각 변형이 주장하는 조건이 실제로
성립했음을 증명한다. (ii) \textbf{결과의 불변식}: 빨라진 코드가
같은 것을 계산하는지 기계적으로 판정해야 한다. 우리는 출력
스파이크의 비트 체크섬(V장)을 모든 변형에 요구한다. (iii)
\textbf{단일 요인 절제}: 최적화를 하나씩 누적한 빌드
변형(A0--A4)을 같은 하네스에서 측정한다.

하네스는 이 세 요소를 위해 설계되었다: 입력은 플래시에 고정된
CIFAR-10 테스트 이미지 8장(SD카드$\cdot$JPEG$\cdot$표시 경로 배제),
인코딩 난수는 시드 고정 xorshift32, 시간은 DWT 사이클
카운터(1사이클 해상도), 워밍업 1장은 제외한다. 하네스 자신도 검증
대상이다: 독립 검토에서 타임베이스 인터럽트 침투(약 0.1\%
노이즈)$\cdot$누산기 자릿수$\cdot$결과 회수 주소 오류 3건이 발견되어
수정 후 절제 변형(A0--A4) 전량을 재측정했다. 선행 연구의 $2{,}407.0$ms는 1ms 해상도
시계$\cdot$표시 루프 포함$\cdot$다른 입력 집합의 수치이므로, 본
하네스로 동일 코드 경로를 재현한 $2{,}328.5$ms(차이 3.3\%)를
기준선(A0)으로 삼고 이후 모든 배속을 하네스 내부 상대 비교로만
계산한다.

\subsection{페치 가설의 반증}\label{sec:falsify}

사전 하드웨어 조사는 페치 병목을 강하게 시사했었다. 이 보드는 옵션
바이트에 의해 플래시의 ITCM 별칭 주소에서 부팅하는데 그 경로의
캐시인 ART 가속기가 꺼져 있었고(\texttt{FLASH\_ACR}$=$0x006),
마이크로벤치마크에서 ART 활성화는 해당 경로를 1.4배 이상, 코드의
ITCM-RAM 이전은 최대 2.5배 가속했다. 그러나 실제 추론 워크로드의
절제 실측(V장 표~\ref{tab:ablation})에서 두 최적화는 \textbf{모두
실질 무효과}였다($\le$0.12\%; A2의 소폭 퇴행은 ITCM 왕복 베니어
오버헤드로 추정된다): ART 활성화(A1)
$2{,}328.5\rightarrow2{,}328.9$ms, 핫 커널 $4{,}392$B의 ITCM
배치(A2) $\rightarrow2{,}331.2$ms.

반증의 기제는 두 가지다. 첫째, \textbf{부팅 주소와 실행 주소는
다르다}. 리셋 벡터는 ITCM 별칭에서 읽히지만 링커는 코드를 AXIM
주소에 배치했으므로, 이후 모든 페치는 AXIM 경로로 나가고 그 경로는
이미 L1 I-캐시가 담당한다. ART는 ITCM 경로 전용이라 개입할 자리가
없다. 둘째, \textbf{핫 명령 집합이 캐시보다 훨씬 작다}. 실행 시간의
대부분을 차지하는 내적 루프는 수백 바이트로 4KB I-캐시에 상주한다.
마이크로벤치마크의 ``큰 커널'' 조건은 6KB 코드를 균일 순회하는 합성
지역성이었고, 실제 워크로드의 지역성과 달랐던 것이다. 이 결과는
마이크로벤치마크 기반 병목 추정의 위험을 보여주는 부정적 증거로서
기록한다.

\subsection{데이터 이동 병목의 확정}\label{sec:datamove}

페치가 아니면 데이터다. dense 커널의 CPL2는 출력 픽셀 2개를 계산할
때마다 퓨전 가중치 전체($147{,}456$B)를 다시 읽는다. 가중치
트래픽만 타임스텝당 약 4.7MB, 이미지당 141.6MB로, 4KB
D-캐시를 상시 축출하는 스트림이다. 실제로 데이터의 \emph{위치}만
바꾼 절제 변형(A3: 가중치$\cdot$LIF 상태$\cdot$스크래치의
DTCM/SRAM 재배치)은 7\%를 얻는 데 그쳤다 --- 144KB 스트림이 4KB
캐시를 축출하는 구조가 그대로인 한 라인필 횟수는 변하지 않으므로,
배치 개선의 여지가 이 수준에 그친다는 방증이다.

총량을 줄이는 유일한 지렛대는 ``곱해 봐야 0인 피연산자를 아예 읽지
않는 것''이고, II장에서 본 것처럼 CPL2 입력의 약 96\%가 그런
피연산자다. 문제는 이로써 명확한 형태를 갖춘다: \emph{희소성을
커널이 소비할 수 있는 형태로 바꾸되, 정수 의미를 비트 단위로
보존하라}. 다음 장이 그 설계다.
